\documentclass[12pt,a4paper,titlepage=false]{scrartcl}

\usepackage{luatex85}
\usepackage{fontspec}
\usepackage[brazilian]{babel}
\usepackage{microtype}
\usepackage{indentfirst}
\usepackage{amsmath,amssymb,mathtools}
\usepackage{graphicx}
\usepackage{geometry}
\usepackage{hyperref}
\usepackage{setspace}
\usepackage{csquotes}
\usepackage{xfrac}
\usepackage[automark]{scrlayer-scrpage}

\setmainfont{Times New Roman}
\setsansfont{Arial}
\setmonofont{Courier New}

\addtokomafont{disposition}{\boldmath}

\clearpairofpagestyles
\rohead{\pagemark}

\onehalfspacing
\recalctypearea
\AfterCalculatingTypearea{\newgeometry{left=3cm, right=2cm, top=3cm, bottom=2cm}}

\title{Variação cambial}
\author{Pedro Schuler \thanks{\href{mailto:pedro.schuler@ufpe.br}{pedro.schuler@ufpe.br}}}
\date{\today}

\begin{document}
\maketitle

\section{O problema}
Algumas vezes podemos nos deparar com cenários em que é necessário efetuar uma multiplicação ao divisão de uma das variáveis independentes. Dessa forma podemos nos perguntar o que acontecerá com o estimador desta variável, com os estimadores de outras variáveis, com a variância desta variável e com o \(\bar{R^{2}}\). Veremos a seguir o impacto de cada uma destas questões.

\section{A solução}

\subsection{Os estimadores}
Utilizando um modelo de regressão simples e assumindo que a variável X será multiplicada por uma constante \(\alpha\).

Sabemos que:
\[\hat{\beta_{1}} = \frac{n\sum{XY} - \sum{X}\sum{Y}}{n\sum{X^{2}}} - (\sum{X})^{2}\]
\[\hat{\beta_{0}} = \bar{Y} - \hat{\beta_{1}}\bar{X}\]

Utilizando as propriedades do somatório e multiplicando X por uma constante \(\alpha\):
\[\hat{\beta_{1}} = \frac{n\sum{\alpha XY} - \sum{\alpha X}\sum{Y}}{n\sum{(\alpha X)^2} - (\sum{\alpha X})^{2}}\]
\[\hat{\beta_{1}} = \frac{\alpha n\sum{XY} - \alpha \sum{X}\sum{Y}}{(\alpha)^2 n\sum{X^{2}} - (\alpha)^{2}(\sum{X})^{2}}\]
\[\hat{\beta_{1}} = \frac{\alpha \times (n\sum{XY} - \sum{X}\sum{Y})}{(\alpha)^2 \times (n\sum{X^{2}} - (\sum{X})^{2})}\]
\[\hat{\beta_{1}} = \frac{\alpha}{\alpha^{2}} \times \frac{n\sum{XY} - \sum{X}\sum{Y})}{n\sum{X^{2}} - (\sum{X})^{2})}\]
\[\hat{\beta_{1}} = \frac{1}{\alpha} \times \frac{n\sum{XY} - \sum{X}\sum{Y})}{n\sum{X^{2}} - (\sum{X})^{2})}\]
\[\hat{\beta_{1}} = \frac{1}{\alpha} \times \hat{\beta_{1}^{\ast}}\]

Onde \(\hat{\beta_{1}^{\ast}}\) é o valor do estimador calculado antes da modificação da variável (multiplicação).

Aplicando a mesma lógica ao \(\beta_{0}\) e sabendo que ao multiplicarmos todos os valores de X por uma constante sua média também será multiplicada pela mesma constante:
\[\hat{\beta_{0}} = \bar{Y} - \hat{\beta_{1}}\bar{X}\]
\[\hat{\beta_{0}} = \bar{Y} - \frac{1}{\alpha}\hat{\beta_{1}^{\ast}}\alpha\bar{X}\]
\[\hat{\beta_{0}} = \bar{Y} - \hat{\beta_{1}^{\ast}}\bar{X}\]
\[\hat{\beta_{0}} = \hat{\beta_{0}^{\ast}}\]

Onde \(\hat{\beta_{1}^{\ast}}\) e \(\hat{\beta_{0}^{\ast}}\) são os valores dos estimadores calculados antes da modificação da variável (multiplicação).

Vemos assim que a variável que sofreu a modificação teve seu estimador multiplicado por \(\sfrac{1}{\alpha}\) enquanto o outro estimador permaneceu inalterado.

\subsection{O \(R^{2}\)}
Sabemos que:
\[R^{2} = 1 - \frac{SQR}{SQT}\]

Sabemos também que:
\[SQT = \sum{(Y_{i} - \bar{Y})^{2}}\]

Como SQT não depende de variações de X, ele permanecerá inalterado.

Procedendo ao cálculo de SQR temos:
\[SQR = \sum{(Y_{i} - \hat{Y_{i}})^{2}}\]
\[SQR = \sum{(Y_{i} - \hat{Y_{i}})^{2}}\]

Resolvendo \(Y_{i} - \hat{Y}_{i}\) temos:
\[Y_{i} - (\hat{\beta_{0}} + \hat{\beta_{1}}X_{i})\]
\[Y_{i} - (\hat{\beta_{0}^{\ast}} + \frac{1}{\alpha} \times \hat{\beta_{1}^{\ast}}\alpha X_{i})\]
\[Y_{i} - (\hat{\beta_{0}^{\ast}} + \hat{\beta_{1}^{\ast}}X_{i})\]

Onde \(\hat{\beta_{1}^{\ast}}\) e \(\hat{\beta_{0}^{\ast}}\) são os valores dos estimadores calculados antes da modificação da variável (multiplicação). Vemos assim que o \(R^{2}\) também permanece inalterado.

\subsection{A variância}
Sabemos que:
\[Var(\hat{\beta_{0}}) = \frac{\sigma^{2}\sum{X_{i}^{2}}}{n\sum{(X_{i} - \bar{X})^{2}}}\]
\[Var(\hat{\beta_{1}}) = \frac{\sigma^{2}}{\sum{(X_{i} - \bar{X})^{2}}}\]

Calculando a \(Var(\hat{\beta_{1}})\) ao multiplicar a variável X por \(\alpha\) temos:
\[Var(\hat{\beta_{1}}) = \frac{\sigma^{2}}{\sum{(\alpha X_{i} - \alpha\bar{X})^{2}}}\]
\[Var(\hat{\beta_{1}}) = \frac{\sigma^{2}}{\sum{((\alpha) \times (X_{i} - \bar{X}))^{2}}}\]
\[Var(\hat{\beta_{1}}) = \frac{\sigma^{2}}{\alpha^{2} \times \sum{(X_{i} - \bar{X})^{2}}}\]
\[Var(\hat{\beta_{1}}) = \frac{1}{\alpha^{2}} \times \frac{\sigma^{2}}{\sum{(X_{i} - \bar{X})^{2}}}\]
\[Var(\hat{\beta_{1}}) = \frac{1}{\alpha^{2}} \times Var(\hat{\beta_{1}^{\ast}})\]

Onde \(Var(\hat{\beta_{1}^{\ast}})\) é o valor da variância de \(\hat{\beta_{1}}\) calculado antes da modificação da variável (multiplicação).

Repetindo o cálculo para \(Var(\hat{\beta_{0}})\):
\[Var(\hat{\beta_{0}}) = \frac{\sigma^{2}\sum{(\alpha X_{i})^{2}}}{n\sum{(\alpha X_{i} - \alpha\bar{X})^{2}}}\]
\[Var(\hat{\beta_{0}}) = \frac{\sigma^{2}\alpha^{2}\sum{X_{i}^{2}}}{n\sum{(\alpha \times (X_{i} - \bar{X}))^{2}}}\]
\[Var(\hat{\beta_{0}}) = \frac{\alpha^{2}\sigma^{2}\sum{X_{i}^{2}}}{\alpha^{2}n\sum{(X_{i} - \bar{X})^{2}}}\]
\[Var(\hat{\beta_{0}}) = \frac{\alpha^{2}}{\alpha^{2}} \times \frac{\sigma^{2}\sum{X_{i}^{2}}}{n\sum{(X_{i} - \bar{X})^{2}}}\]
\[Var(\hat{\beta_{0}}) = 1 \times \frac{\sigma^{2}\sum{X_{i}^{2}}}{n\sum{(X_{i} - \bar{X})^{2}}}\]
\[Var(\hat{\beta_{0}}) = \frac{\sigma^{2}\sum{X_{i}^{2}}}{n\sum{(X_{i} - \bar{X})^{2}}}\]
\[Var(\hat{\beta_{0}}) = Var(\hat{\beta_{0}^{\ast}})\]

Onde \(Var(\hat{\beta_{0}^{\ast}})\) é o valor da variância de \(\hat{\beta_{0}}\) calculado antes da modificação da variável (multiplicação).

Ou seja: a variável que sofreu a modificação terá sua variância multiplicada por \(\sfrac{1}{\alpha^{2}}\) e a variância das outras variáveis permanecerão inalteradas.
\end{document}
